\label{ch:modular-pva-quantization}

\begin{abstract}
This note reorganizes the entire discussion of the present thread into a single fortified
package, written from first principles and separated carefully into three layers:

\begin{enumerate}[label=\textnormal{(\roman*)},leftmargin=2.6em]
\item \emph{external input} from the literature;
\item \emph{formal algebraic results} proved internally in this note once a modular bar datum is given;
\item \emph{conditional comparison statements} that connect the algebraic package to the
three-dimensional holomorphic-topological Poisson sigma model of Khan--Zeng.
\end{enumerate}

The central principle is that \emph{trees govern chirality and loops govern modularity}.
Disk-level quadratic duality controls only genus zero. A genuinely modular deformation
theory must already contain a primitive genus-raising one-edge operator. We make this
precise by constructing a stable-graph modular bar coalgebra, defining the complete
filtered deformation dg Lie algebra
\[
\Defmod(A_{\cl}(V)) := \Coder(\Bmod(\A))[-1],
\]
and proving the universal obstruction recursion for lifting a genus-zero Maurer--Cartan
class to all genera. We then pass to a local variational coordinate model for filtered
quadratic Poisson vertex algebras, identify the first genus-raising operator with a reduced
odd Laplacian
\[
D_1=\Delta_{\mathrm{cyc}},
\]
construct the associated variational modular class, prove a vanishing theorem for a broad
triangular $W$-normal form, and compute explicitly that the first modular obstruction
vanishes for classical $W_3$.

Finally, in the local, translation-invariant, filtration-preserving, reduced conformal-weight
zero sector, we compute the relevant shifted $H^1$ for $W_3$ and prove that the only
nontrivial genus-one lift direction is the central-parameter direction $\partial_c P_c$.
Thus, in the relevant one-loop sector, every genus-one lift is gauge equivalent to a
renormalization of the central parameter.

The resulting picture gives a precise candidate for the missing modular homotopy theory
behind deformation quantization of Poisson vertex algebras to vertex algebras, while
remaining honest about which inputs are presently formal, conditional, or conjectural.
\end{abstract}

\tableofcontents

\section{Status conventions and scope}

The note is intentionally stratified.

\begin{convention}[Three levels of assertion]
Throughout the note we distinguish three kinds of statements.
\begin{enumerate}[leftmargin=2.4em]
\item \textbf{External input.} These are statements imported from the literature:
Khan--Zeng on the three-dimensional holomorphic-topological Poisson sigma model,
Gui--Li--Zeng on quadratic duality for chiral algebras, Francis--Gaitsgory on chiral
Koszul duality, Costello--Dimofte--Gaiotto on boundary chiral algebras in 3d
holomorphic twists, Gaiotto--Kulp--Wu on higher operations, and related works
\cite{KZ25,GLZ22,FG12,CDG20,GKW25}.
\item \textbf{Formal algebraic results.} These are theorems proved in this note once one
has fixed a modular bar datum. They include the stable-graph bar construction, the
filtered deformation dg Lie algebra, the universal obstruction recursion, the coordinate
formula for the genus-raising operator, the variational modular class, the triangular
$W$-normal-form vanishing theorem, and the $W_3$ relevant cohomology computation.
\item \textbf{Conditional comparison results.} These are statements linking the formal
algebra to the Poisson sigma model. They depend on analytic and geometric inputs
not proved here, such as the existence of the appropriate resolved chiral/factorization
datum attached to a given PVA and the existence of half-space perturbative quantization.
\end{enumerate}
\end{convention}

\begin{remark}[Why this stratification matters]
The guiding philosophy of the thread was to push the frontier without pretending to
have already solved every analytic problem. In particular, the algebraic stable-graph
package below is self-contained once a modular bar datum is fixed, but the extraction of
that datum from a general PVA via compactified configuration spaces and the identification
with the three-dimensional Poisson sigma model remain conditional at present.
\end{remark}

\section{External input from the literature}

We record the precise ingredients from the literature that motivate and support the
new constructions.

\subsection{Poisson vertex algebras and the 3d holomorphic-topological Poisson sigma model}

Khan and Zeng introduce a mixed holomorphic-topological gauge theory in three dimensions
attached to a freely generated Poisson vertex algebra (PVA), show that gauge invariance
holds if and only if the $\lambda$-bracket Jacobi identity holds, prove that the presence
of a Virasoro element upgrades the symmetry from holomorphic-topological to fully
topological, and propose the resulting three-dimensional Poisson sigma model as a route
toward deformation quantization of PVAs to vertex algebras \cite{KZ25}. Their
boundary picture on the half-space $\mathbb R_+\times \mathbb C$ is the physical source of
the genus-zero quantization problem considered here.

\subsection{Quadratic duality for chiral algebras}

Gui, Li, and Zeng define quadratic duality for chiral algebras, prove an injective map
\[
\Hom(A,B)\hookrightarrow \MC(A^!\otimes B),
\]
show bijectivity under effectiveness hypotheses in degree zero, and in the inhomogeneous
quadratic setting introduce twisted pairs $(B,B^\circ,S)$ together with the curved
Maurer--Cartan equation
\[
\mu\big((S+\alpha)\boxtimes (S+\alpha)\big)=0.
\]
They also state explicitly that a full homotopy/Koszul treatment, parallel to the
associative story, is still missing \cite{GLZ22}. This gap is one of the main
algebraic motivations for the present note.

\subsection{Chiral/factorization homotopy theory}

Francis and Gaitsgory place chiral and factorization structures in a common homotopy-theoretic
framework and recast their relation as a form of chiral Koszul duality
\cite{FG12}. This supplies the correct derived ambient category in which
one should expect a modular refinement of disk-level chiral duality to live.

\subsection{Boundary algebras, higher operations, and line operators}

Costello, Dimofte, and Gaiotto show that in 3d holomorphic twists bulk local operators form
a commutative chiral algebra with shifted Poisson bracket, while boundary operator algebras
naturally carry higher $A_\infty$ structure \cite{CDG20}. Gaiotto,
Kulp, and Wu construct higher operations in holomorphic-topological perturbation theory,
prove the relevant quadratic consistency identities, and interpret them as a Wess--Zumino
type consistency for BRST symmetry \cite{GKW25}. Dimofte, Niu, and Py show
that line operators in perturbative 3d holomorphic QFT are represented by modules for a
Koszul-dual $A_\infty$-algebra $A^!$ and are organized by Maurer--Cartan data with an
$A_\infty$ Yang--Baxter flavor \cite{DNP25}. All of this reinforces the thesis
that \emph{a single Maurer--Cartan object should package higher operations, anomalies,
and OPE compatibilities}. 

\subsection{Analogies with deformation quantization}

Cattaneo and Felder interpret Kontsevich's deformation quantization via a topological
sigma model on the disk \cite{CF00}. Tamarkin proves that the deformation
complex of an $E_d$-algebra carries an $E_{d+1}$-structure \cite{Tam02}. These
results suggest that the correct receptacle for quantization of PVAs should not be merely
a disk-level dg Lie algebra, but a higher and genuinely modular deformation object.

\section{From first principles: genus zero, loops, and the one-edge principle}

The deepest structural lesson of the thread can be stated in one sentence:

\begin{center}
\emph{A deformation theory that only remembers trees cannot see modular obstructions.}
\end{center}

We make that precise.

\subsection{Filtered quadratic PVAs}

\begin{definition}[Filtered quadratic PVA]
A \emph{filtered quadratic Poisson vertex algebra} is a differential commutative algebra
freely generated by fields $u^1,\dots,u^N$ and their derivatives, together with a
$\lambda$-bracket encoded by a skew-adjoint matrix differential operator
\[
\Pi^{ab}(\partial)=\sum_{m\ge 0}\Pi^{ab}_m\,\partial^m,
\]
whose coefficients are filtered polynomials of total generator degree at most $2$.
\end{definition}

The classical $W$-algebras discussed by Khan--Zeng fit this paradigm, including the
classical $W_3$ example \cite{KZ25}.

\subsection{Resolved classical datum}

\begin{assumption}[Resolved classical chiral/factorization datum]
Let $X$ be a smooth complex curve and let $V$ be a filtered quadratic PVA on $X$.
Assume there exists a cofibrant resolved classical chiral/factorization object
\[
\A=\widetilde A_{\cl}(V)
\]
whose genus-zero shadow is $V$, and which is dualizable, finite type, conilpotent,
and complete in the sense made precise below.
\end{assumption}

\begin{remark}
This assumption isolates the step that is currently external to the purely algebraic theory.
The remainder of the formal constructions depend only on the existence of an input
modular bar datum built from $\A$.
\end{remark}

\subsection{The one-edge principle}

The naive temptation is to take the disk-level deformation complex and merely add a
genus grading. That fails.

\begin{proposition}[Necessity of a genus-raising primitive; \ClaimStatusProvedHere]
Suppose a filtered deformation dg Lie algebra $L$ is equipped with a complete genus
filtration $F^\bullet$ and its differential preserves genus exactly. Then any genus-zero
Maurer--Cartan element $\Theta_0\in \MC(\gr_F^0 L)$ extends tautologically to a full
Maurer--Cartan element of $L$ without further obstruction.
\end{proposition}

\begin{proof}
If the differential preserves genus and the Lie bracket adds genera, then the Maurer--Cartan
equation decomposes strictly by genus. A solution in genus zero has no forcing term in
positive genus, so the same element, viewed inside the completed filtered algebra, already
solves the full equation.
\end{proof}

\begin{remark}[Trees versus loops]
Therefore any nontrivial modular obstruction theory must contain, at the chain level,
a primitive operator that \emph{raises genus}. In the stable-graph bar construction below,
that primitive is the nonseparating one-edge expansion operator. This is the algebraic
counterpart of the first loop correction in perturbation theory.
\end{remark}

\section{Stable-graph modular bar data and the modular bar coalgebra}

This section is the formal heart of the new algebraic package.

\subsection{Stable graphs}

\begin{definition}[Connected stable graph]
A \emph{connected stable graph} $\Gamma$ consists of the following data:
\begin{enumerate}[leftmargin=2.4em]
\item a finite set $V(\Gamma)$ of vertices;
\item a finite set $\mathrm{Fl}(\Gamma)$ of flags;
\item a partition $\mathrm{Fl}(\Gamma)=\bigsqcup_{v\in V(\Gamma)}\mathrm{Fl}(v)$;
\item an involution on $\mathrm{Fl}(\Gamma)$, whose $2$-cycles are internal edges and
whose fixed points are external legs;
\item a genus label $g(v)\in \mathbb Z_{\ge 0}$ for each vertex,
\end{enumerate}
subject to the stability condition
\[
2g(v)-2+|\mathrm{Fl}(v)|>0
\qquad\text{for all }v\in V(\Gamma).
\]
Its total genus is
\[
g(\Gamma)=b_1(\Gamma)+\sum_{v\in V(\Gamma)}g(v).
\]
\end{definition}

Write $E(\Gamma)$ for the set of internal edges and
\[
\orline(\Gamma):=\det(kE(\Gamma))\otimes \det(H_1(\Gamma;k))^{-1}
\]
for the orientation line.

\subsection{One-edge expansions}

Let $C$ be a reduced complete finite-type graded collection on finite sets:
\[
F\longmapsto C(F),
\]
functorial for bijections of finite sets.

\begin{definition}[One-edge expansion set]
For a finite set $F$, let $\OE(F)$ denote the set of connected stable graphs whose
external leg set is $F$ and which have exactly one internal edge. Such a graph is either:
\begin{enumerate}[leftmargin=2.4em]
\item \emph{separating}: two vertices joined by one edge;
\item \emph{nonseparating}: one vertex carrying one loop.
\end{enumerate}
\end{definition}

\begin{definition}[Modular bar datum]
A \emph{modular bar datum} on $C$ consists of:
\begin{enumerate}[leftmargin=2.4em]
\item an internal differential
\[
d:C(F)\to C(F)
\]
of degree $+1$ for each finite set $F$;
\item for each $\epsilon\in \OE(F)$, a degree $+1$ map
\[
\Delta_\epsilon:
s\,C(F)\longrightarrow
\orline(\epsilon)\otimes \bigotimes_{u\in V(\epsilon)} s\,C(\mathrm{Fl}_\epsilon(u));
\]
\item the \emph{codimension-two cancellation axiom}: for every connected stable graph
$T$ with exactly two internal edges, the signed sum of the two ways of producing $T$
by successive one-edge expansions vanishes;
\item the anticommutation relation that $d$ anticommutes with every $\Delta_\epsilon$.
\end{enumerate}
\end{definition}

\begin{remark}[Why this is the right primitive datum]
The usual bar construction is governed by codimension-one collisions on configuration
spaces. In the modular setting, codimension-one boundary strata correspond to one-edge
degenerations of stable graphs. Thus the primitive input is not a generic higher operator
but precisely the family of one-edge expansion maps.
\end{remark}

\subsection{The complete modular bar coalgebra}

For a connected stable graph $\Gamma$, define
\[
C[\Gamma]:=
\orline(\Gamma)\otimes\bigotimes_{v\in V(\Gamma)} s\,C(\mathrm{Fl}(v)).
\]

\begin{definition}[Complete modular bar object]
The \emph{complete modular bar object} of a modular bar datum $C$ is
\[
\Bmod(C):=
\prod_{[\Gamma]\in \StGraph^{\mathrm{conn}}}
\big(C[\Gamma]\big)_{\mathrm{Aut}(\Gamma)},
\]
completed with respect to total genus.
\end{definition}

\begin{definition}[Internal and expansion differentials]
Define degree-$+1$ endomorphisms
\[
D_{\mathrm{int}},\ D_{\mathrm{exp}}:\Bmod(C)\to \Bmod(C)
\]
on graph generators as follows.
\begin{enumerate}[leftmargin=2.4em]
\item $D_{\mathrm{int}}$ applies the internal differential $d$ to one vertex label and
sums over vertices.
\item $D_{\mathrm{exp}}$ replaces a chosen vertex $v$ by a one-edge expansion
$\epsilon\in \OE(\mathrm{Fl}(v))$ and applies $\Delta_\epsilon$ to the label at $v$,
with the sign determined by the orientation line and Koszul rule.
\end{enumerate}
Set
\[
D:=D_{\mathrm{int}}+D_{\mathrm{exp}}.
\]
\end{definition}

\begin{theorem}[Stable-graph modular bar construction; \ClaimStatusProvedHere]\label{thm:modular-bar}
Let $C$ be a modular bar datum. Then:
\begin{enumerate}[leftmargin=2.4em]
\item $D^2=0$;
\item $\Bmod(C)$ is a complete conilpotent modular dg coalgebra;
\item if $F^g\Bmod(C)$ denotes the subspace spanned by graphs of total genus $\ge g$, then
\[
\Bmod(C)\cong \varprojlim_g \Bmod(C)/F^g\Bmod(C),
\]
and this filtration is exhaustive and complete.
\end{enumerate}
\end{theorem}

\begin{proof}
The square of $D$ has three types of terms.

\smallskip
\noindent
\emph{(i) Two internal differentials.}
These vanish because $d^2=0$ on every vertex label.

\smallskip
\noindent
\emph{(ii) One internal differential and one expansion.}
These cancel because $d$ anticommutes with every one-edge expansion map $\Delta_\epsilon$.

\smallskip
\noindent
\emph{(iii) Two one-edge expansions.}
These are exactly the codimension-two stable-graph boundary terms. Every graph with two
internal edges can be obtained by two successive one-edge expansions, and the codimension-two
cancellation axiom says the signed sum of the two resulting paths is zero.

Therefore $D^2=0$.

For conilpotence, use the reduced coproduct defined by extracting a nontrivial connected
stable subgraph. Repeated reduced coproduct strictly decreases the number of internal
edges, hence eventually vanishes on every graph class. Completeness is immediate from
the completed product over total genus.
\end{proof}

\begin{remark}[Tree-level truncation]
The genus-zero subquotient of $\Bmod(C)$ is obtained by restricting to connected stable
graphs of total genus zero, i.e.\ trees all of whose vertices have genus zero. We denote
that tree-level coalgebra by $\Bch(C)$.
\end{remark}

\subsection{Separating and nonseparating pieces}

Split the expansion operator into
\[
D_{\mathrm{exp}}=D_{\mathrm{sep}}+D_{\mathrm{nsep}},
\]
where $D_{\mathrm{sep}}$ sums over separating one-edge expansions and $D_{\mathrm{nsep}}$
over nonseparating one-edge expansions. Set
\[
D_0:=D_{\mathrm{int}}+D_{\mathrm{sep}},
\qquad
D_1:=D_{\mathrm{nsep}}.
\]

\begin{proposition}[The one-edge decomposition; \ClaimStatusProvedHere]\label{prop:D0D1}
With the notation above:
\begin{enumerate}[leftmargin=2.4em]
\item $D_0$ preserves total genus;
\item $D_1$ raises total genus by one;
\item $D_0^2=0$, $D_0D_1+D_1D_0=0$, and $D_1^2=0$.
\end{enumerate}
\end{proposition}

\begin{proof}
Separating one-edge expansions preserve the first Betti number and total genus, while
nonseparating one-edge expansions increase the first Betti number by one. The identities
follow by taking the genus-homogeneous pieces of the equality $(D_0+D_1)^2=0$.
\end{proof}

\subsection{Curved extension via $S$-tails}

Gui--Li--Zeng show that in the inhomogeneous quadratic setting the correct input is not
merely a CDG algebra, but a twisted pair $(B,B^\circ,S)$, and that the passage to the
associated CDG algebra forgets information \cite{GLZ22}. The modular extension
should therefore remember the element $S$ geometrically.

\begin{definition}[Curved one-tail expansion]
Given a twisted pair $(B,B^\circ,S)$, an \emph{$S$-tail expansion} is the operation that
attaches a new univalent $S$-decorated vertex to a pre-existing corolla. Formally, it is
a unary expansion map
\[
\Delta_S:s\,C(F)\to s\,C(F\sqcup \{\bullet,\star\})
\]
followed by contraction of the new internal edge, with the new univalent vertex labeled by $S$.
\end{definition}

\begin{remark}
The correct curved modular bar object is obtained by adjoining the $S$-tail expansions
to the set of primitive degenerations. Provided they satisfy the same codimension-two
compatibilities with the one-edge expansions, the proof of \cref{thm:modular-bar} goes
through verbatim. This is the modular analogue of the twisted Maurer--Cartan equation
\[
\mu\big((S+\alpha)\boxtimes (S+\alpha)\big)=0
\]
of Gui--Li--Zeng.
\end{remark}

\section{The modular deformation dg Lie algebra}

We now pass from the coalgebra to the deformation complex.

\begin{definition}[Modular deformation algebra]
Let $C$ be a modular bar datum. Define
\[
L_{\mathrm{mod}}(C):=\Coder(\Bmod(C))[-1].
\]
If $\A$ is a resolved classical chiral/factorization datum with underlying modular bar datum
$C(\A)$, write
\[
\Defmod(A_{\cl}(V)):=L_{\mathrm{mod}}(C(\A)).
\]
\end{definition}

\begin{proposition}[Convolution description; \ClaimStatusProvedHere]
A coderivation of $\Bmod(C)$ is determined by its projection to corollas. Hence there is
a canonical identification
\[
L_{\mathrm{mod}}(C)\cong \Hom(\Bmod(C),sC)[-1]
\]
with dg Lie bracket given by commutator.
\end{proposition}

\begin{proof}
This is the standard universal property of cofree conilpotent coalgebras, applied in the
completed stable-graph setting.
\end{proof}

\subsection{The genus filtration}

\begin{definition}[Genus-raise filtration]
For $\delta\in L_{\mathrm{mod}}(C)$ define
\[
\delta\in F^pL_{\mathrm{mod}}(C)
\iff
\delta(F^g\Bmod(C))\subseteq F^{g+p}\Bmod(C)
\quad\text{for all }g\ge 0.
\]
\end{definition}

\begin{theorem}[Formal structure theorem; \ClaimStatusProvedHere]\label{thm:formal-structure}
Let $C$ be a modular bar datum. Then:
\begin{enumerate}[leftmargin=2.4em]
\item $L_{\mathrm{mod}}(C)$ is a complete filtered dg Lie algebra with differential
\[
d_L(\delta):=[D,\delta];
\]
\item the filtration is Lie-multiplicative:
\[
[F^pL_{\mathrm{mod}},F^qL_{\mathrm{mod}}]\subseteq F^{p+q}L_{\mathrm{mod}};
\]
\item the positive-genus ideal $F^1L_{\mathrm{mod}}$ is pronilpotent;
\item the genus-zero associated graded piece identifies with the ordinary tree-level deformation complex:
\[
\gr_F^0 L_{\mathrm{mod}}(C)\cong \Coder(\Bch(C))[-1].
\]
\end{enumerate}
\end{theorem}

\begin{proof}
Coderivations of any dg coalgebra form a dg Lie algebra under commutator. Completeness
is inherited from the completeness of $\Bmod(C)$. If $\delta$ raises genus by at least $p$
and $\eta$ by at least $q$, then both compositions $\delta\eta$ and $\eta\delta$ raise genus
by at least $p+q$, proving Lie-multiplicativity. Pronilpotence of $F^1$ follows because on
the finite quotient by $F^N$, sufficiently long brackets vanish.

For the associated graded in genus zero, only total-genus-zero graphs survive, i.e.\ trees
with vertex genera all zero. The nonseparating operator $D_1$ dies in the quotient, so the
differential reduces to the tree-level bar differential.
\end{proof}

\begin{corollary}[No modular theory without loops; \ClaimStatusProvedHere]
If $D_1=0$, then every genus-zero Maurer--Cartan class extends uniquely to a full
Maurer--Cartan class. In particular, the nonseparating one-edge piece is the minimal
algebraic source of genuine modular obstructions.
\end{corollary}

\begin{proof}
Apply \cref{prop:D0D1} and the argument of the proposition from \S3.
\end{proof}

\subsection{Genus-zero comparison with quadratic chiral duality}

\begin{proposition}[Comparison with Gui--Li--Zeng; \ClaimStatusProvedHere]
Suppose the genus-zero truncation of $C(\A)$ is the ordinary quadratic chiral bar datum
of a dualizable quadratic chiral algebra $A_{\cl}(V)$. Then the genus-zero Maurer--Cartan
problem in $\gr_F^0\Defmod(A_{\cl}(V))$ recovers the Gui--Li--Zeng Maurer--Cartan problem:
\[
\Hom(A,B)\hookrightarrow \MC(A^!\otimes B)
\]
in the strictly quadratic case, and the twisted equation
\[
\mu\big((S+\alpha)\boxtimes (S+\alpha)\big)=0
\]
in the inhomogeneous quadratic case.
\end{proposition}

\begin{proof}
On the genus-zero quotient only trees survive, so the modular bar coalgebra reduces to the
ordinary tree-level chiral bar coalgebra. The corresponding coderivation dg Lie algebra is
precisely the disk-level deformation complex. The Gui--Li--Zeng identification is therefore
the genus-zero shadow of the present modular theory.
\end{proof}

\section{Maurer--Cartan obstruction theory}

We now extract the exact obstruction groups.

\subsection{Formal obstruction recursion}

Let $L=L_{\mathrm{mod}}(C)$ and write $d_0=[D_0,-]$, $d_1=[D_1,-]$.

\begin{theorem}[Universal genus-one obstruction; \ClaimStatusProvedHere]\label{thm:Ob1}
Let $\Theta_0\in \MC(\gr_F^0L)$. Then the class
\[
\Ob_1(\Theta_0):=[d_1\Theta_0]
\in
H^2\big((\gr_F^1L)^{\Theta_0},d_0^{\Theta_0}\big),
\qquad
d_0^{\Theta_0}(x):=d_0x+[\Theta_0,x],
\]
is well-defined. Moreover,
\[
\Ob_1(\Theta_0)=0
\iff
\exists\,\Theta_1\in F^1L/F^2L
\text{ such that }
\Theta_0+\Theta_1
\]
solves the Maurer--Cartan equation modulo $F^2L$.
\end{theorem}

\begin{proof}
Expand the Maurer--Cartan equation for $\Theta=\Theta_0+\Theta_1$ modulo $F^2$:
\[
0=d\Theta+\frac12[\Theta,\Theta]
=
\Big(d_0\Theta_0+\frac12[\Theta_0,\Theta_0]\Big)
+
\Big(d_0\Theta_1+d_1\Theta_0+[\Theta_0,\Theta_1]\Big).
\]
The first parenthesis vanishes because $\Theta_0$ is genus-zero Maurer--Cartan, hence the
lifting equation is
\[
d_0^{\Theta_0}\Theta_1=-d_1\Theta_0.
\]
Thus the lift exists iff the cohomology class of $d_1\Theta_0$ vanishes.

To prove that $d_1\Theta_0$ is closed, use \cref{prop:D0D1}. Since $D_0^2=0$,
$D_0D_1+D_1D_0=0$, and $D_1^2=0$, we obtain
\[
d_0^2=0,\qquad d_0d_1+d_1d_0=0.
\]
Applying this to the genus-zero Maurer--Cartan equation
\[
d_0\Theta_0+\frac12[\Theta_0,\Theta_0]=0
\]
and using that $d_1$ is a derivation of the Lie bracket yields
\[
d_0^{\Theta_0}(d_1\Theta_0)=0.
\]
\end{proof}

\begin{theorem}[Higher-genus recursion; \ClaimStatusProvedHere]\label{thm:Obg}
Suppose
\[
\Theta_{<g}:=\Theta_0+\Theta_1+\cdots+\Theta_{g-1}
\]
solves the Maurer--Cartan equation modulo $F^gL$. Then the genus-$g$ forcing term
\[
R_g(\Theta_{<g})
:=
d_1\Theta_{g-1}
+\frac12\sum_{i=1}^{g-1}[\Theta_i,\Theta_{g-i}]
\]
is $d_0^{\Theta_0}$-closed and determines a canonical obstruction class
\[
\Ob_g(\Theta_{<g})
:=
[R_g(\Theta_{<g})]
\in
H^2\big((\gr_F^gL)^{\Theta_0},d_0^{\Theta_0}\big).
\]
It vanishes if and only if there exists $\Theta_g\in F^gL/F^{g+1}L$ such that
\[
\Theta_{\le g}:=\Theta_0+\cdots+\Theta_g
\]
solves the Maurer--Cartan equation modulo $F^{g+1}L$.
\end{theorem}

\begin{proof}
Take the genus-$g$ part of the full Maurer--Cartan equation. Since $D_1$ raises genus by
exactly one, the only differential terms are $d_0\Theta_g$ and $d_1\Theta_{g-1}$, while the
quadratic part is the sum over pairs $(i,g-i)$. Thus
\[
d_0^{\Theta_0}\Theta_g=-R_g(\Theta_{<g}).
\]
The same computation as in \cref{thm:Ob1} shows that $R_g$ is closed. The vanishing
criterion is immediate.
\end{proof}

\begin{corollary}[Full modular lift criterion; \ClaimStatusProvedHere]
A full modular Maurer--Cartan lift
\[
\Theta=\Theta_0+\Theta_1+\Theta_2+\cdots\in \MC(L)
\]
exists if and only if
\[
\Ob_g(\Theta_{<g})=0
\qquad\text{for all }g\ge 1.
\]
\end{corollary}

\subsection{Conditional comparison with the Khan--Zeng half-space theory}

We now formulate the theorem package in the form best suited for later application to the
three-dimensional Poisson sigma model.

\begin{assumption}[Half-space tree-level quantization]
Assume the Khan--Zeng half-space theory on $\mathbb R_+\times X$ with Neumann boundary
condition admits renormalized perturbation theory through tree level.
\end{assumption}

\begin{theorem}[Conditional tree-level boundary quantization; \ClaimStatusConditional]\label{thm:conditional-tree}
Let $V$ be a filtered quadratic PVA and let $\A=\widetilde A_{\cl}(V)$ satisfy the resolved
classical datum hypotheses. Under the half-space tree-level quantization assumption there
exist:
\begin{enumerate}[leftmargin=2.4em]
\item a genus-zero Maurer--Cartan class
\[
\Theta_0\in \MC(\gr_F^0\Defmod(A_{\cl}(V)));
\]
\item a genus-zero boundary chiral/factorization algebra $Q_0(V)$,
\end{enumerate}
such that:
\begin{enumerate}[label=\textnormal{(\alph*)},leftmargin=2.8em]
\item the genus-zero shadow of $\Theta_0$ is the classical $\lambda$-bracket of $V$;
\item the classical limit of $Q_0(V)$ is $A_{\cl}(V)$;
\item under the genus-zero comparison with Gui--Li--Zeng, $\Theta_0$ is the disk-level
Maurer--Cartan class encoding the boundary quantization map
\[
A_{\cl}(V)\longrightarrow Q_0(V).
\]
\end{enumerate}
If $V$ contains a Virasoro element satisfying the Khan--Zeng hypotheses, then $\Theta_0$
lies in the topological subcomplex
\[
\Deftop(A_{\cl}(V))\subseteq \Defmod(A_{\cl}(V)).
\]
\end{theorem}

\begin{remark}
This theorem is conditional because the formal algebraic package constructed here does not
by itself prove the existence of the half-space perturbative quantization. What it does is
identify the \emph{exact deformation-theoretic object} into which that boundary problem
must land.
\end{remark}

\begin{corollary}[Conditional modular quantization criterion; \ClaimStatusConditional]
Under the hypotheses of \cref{thm:conditional-tree}, if all obstruction classes vanish,
then there exists
\[
\Theta\in \MC(\Defmod(A_{\cl}(V)))
\]
whose genus-zero shadow is the classical $\lambda$-bracket and whose boundary value is the
desired quantized chiral algebra. Under the Virasoro hypotheses the lift can be chosen in
$\Deftop(A_{\cl}(V))$.
\end{corollary}

\section{Coordinate realization on local variational polyvectors}

We now construct the concrete coordinate model for the first genus-raising operator.

\subsection{Local polyvectors}

Let
\[
\widehat{\cV}:=
\widehat{\mathbb C}[u^a_{(n)}\mid a=1,\dots,N,\ n\ge 0]
\]
be the completed differential polynomial algebra, with $\partial u^a_{(n)}=u^a_{(n+1)}$.
Introduce odd variables $\theta_a^{(n)}$ satisfying
\[
\partial \theta_a^{(n)}=\theta_a^{(n+1)}.
\]

\begin{definition}[Local polyvectors]
The space of local polyvectors is
\[
\Xloc^\bullet(\widehat{\cV})
:=
\widehat{\cV}[\theta_a^{(n)}]\big/\partial\big(\widehat{\cV}[\theta_a^{(n)}]\big).
\]
We write $\int F$ for the class of $F$ modulo total derivatives. The usual variational
Schouten bracket on $\Xloc^\bullet$ will be denoted by $[\ ,\ ]$.
\end{definition}

\begin{remark}[Shift conventions]
The ordinary Poisson cohomological degree is the polyvector degree. Thus vector fields are
degree $1$ and bivectors are degree $2$. The shifted Maurer--Cartan deformation complex is
obtained from $\Xloc^{\bullet+1}$. In this note we freely pass between the two conventions,
always indicating which degree is meant.
\end{remark}

Every local bivector has a reduced representative of the form
\[
P=\frac12\int \theta_a\,A^{ab}(\partial)\,\theta_b
\]
with $A^{ab}(\partial)$ skew-adjoint and all derivatives moved onto the right-hand factor.

For a filtered quadratic PVA operator
\[
\Pi^{ab}(\partial)=\sum_{m\ge 0}\Pi_m^{ab}\,\partial^m,
\]
write
\[
P_\Pi:=\frac12\int \theta_a\,\Pi^{ab}(\partial)\,\theta_b.
\]
The PVA Jacobi identity is equivalent to the variational master equation
\[
[P_\Pi,P_\Pi]=0.
\]

\subsection{The reduced odd Laplacian}

\begin{definition}[Reduced cyclic odd Laplacian]
Define
\[
\Delta_{\mathrm{cyc}}\!\left(\int F\right)
:=
\sum_{a=1}^N\sum_{n\ge 0}
\int
(-\partial)^n
\frac{\partial^2 F}{\partial u^a_{(n)}\,\partial \theta_a}.
\]
\end{definition}

\begin{remark}[Interpretation]
The operator $\Delta_{\mathrm{cyc}}$ is the coordinate model of the nonseparating one-edge
expansion. Geometrically it contracts the two half-edges of a one-loop self-gluing; in BV
language it is the reduced divergence of a local bivector.
\end{remark}

\begin{theorem}[Coordinate formula for the genus-raising operator; \ClaimStatusProvedHere]\label{thm:D1formula}
Let $\Pi=(\Pi^{ab}(\partial))$ be a filtered quadratic PVA operator. Then
\[
D_1(P_\Pi)=\Delta_{\mathrm{cyc}}P_\Pi
=
\int \sum_{b=1}^N \mathfrak M_\Pi^b(\partial)(\theta_b),
\]
where
\[
\mathfrak M_\Pi^b(\partial)
:=
\sum_{a=1}^N \frac{\delta \Pi^{ab}(\partial)}{\delta u^a},
\qquad
\frac{\delta \Pi^{ab}(\partial)}{\delta u^a}
:=
\sum_{n\ge 0}(-\partial)^n\circ
\frac{\partial \Pi^{ab}(\partial)}{\partial u^a_{(n)}}.
\]
We call $\mathfrak M_\Pi=(\mathfrak M_\Pi^b)_b$ the \emph{variational modular operator}
of $\Pi$.
\end{theorem}

\begin{proof}
Differentiate
\[
P_\Pi=\frac12\int \theta_a\Pi^{ab}(\partial)\theta_b
\]
with respect to $\theta_a$ to contract one odd half-edge, then take the variational
derivative with respect to the even jet variables $u^a_{(n)}$. Each integration by parts
contributes a factor $(-\partial)^n$. Summing over $a$ and $n$ produces exactly the
displayed formula.
\end{proof}

\begin{lemma}[BV properties of $\Delta_{\mathrm{cyc}}$; \ClaimStatusProvedHere]\label{lem:BV}
The operator $\Delta_{\mathrm{cyc}}$ is odd, second-order, and satisfies:
\begin{align*}
\Delta_{\mathrm{cyc}}^2&=0,\\
\Delta_{\mathrm{cyc}}[X,Y]
&=
[\Delta_{\mathrm{cyc}}X,Y]
+(-1)^{|X|+1}[X,\Delta_{\mathrm{cyc}}Y].
\end{align*}
\end{lemma}

\begin{proof}
This is the standard BV-coordinate computation in jet coordinates. The Euler-operator
factors $(-\partial)^n$ already incorporate all integration-by-parts signs. Super-commutativity
forces the square to vanish, and the second-order nature yields the derivation rule for the
Schouten bracket.
\end{proof}

\subsection{The variational modular class}

\begin{theorem}[Variational modular class; \ClaimStatusProvedHere]\label{thm:modclass}
Let $\Pi$ be a filtered quadratic PVA operator. Then
\[
[P_\Pi,\Delta_{\mathrm{cyc}}P_\Pi]=0.
\]
Hence $\Delta_{\mathrm{cyc}}P_\Pi$ determines a canonical local Poisson cohomology class
\[
\Mod(V):=[\Delta_{\mathrm{cyc}}P_\Pi]\in
H^1_{\delta_\Pi}\!\big(\Xloc^\bullet(\widehat{\cV})\big),
\qquad
\delta_\Pi:=[P_\Pi,-].
\]
We call $\Mod(V)$ the \emph{variational modular class} of the filtered quadratic PVA.
\end{theorem}

\begin{proof}
Apply \cref{lem:BV} to the identity $[P_\Pi,P_\Pi]=0$:
\[
0=\Delta_{\mathrm{cyc}}[P_\Pi,P_\Pi]=2[P_\Pi,\Delta_{\mathrm{cyc}}P_\Pi].
\]
Thus $\Delta_{\mathrm{cyc}}P_\Pi$ is $\delta_\Pi$-closed.
\end{proof}

\begin{remark}[Obstruction shadow]
The class $\Mod(V)$ lives in ordinary Poisson cohomology. Under the one-edge comparison
between the nonseparating bar operator and the local variational model, it is the coordinate
shadow of the genus-one forcing term $d_1\Theta_0$ in \cref{thm:Ob1}. In particular, if
$\Delta_{\mathrm{cyc}}P_\Pi=0$ then the coordinate shadow of the first modular obstruction
vanishes.
\end{remark}

\begin{lemma}[Gauge-independence up to exact terms; \ClaimStatusProvedHere]\label{lem:gaugeind}
Let $f$ be an even local functional and define
\[
\Delta_f:=e^{-f}\Delta_{\mathrm{cyc}}e^f.
\]
Then
\[
\Delta_fP_\Pi-\Delta_{\mathrm{cyc}}P_\Pi=[P_\Pi,f].
\]
Hence the cohomology class $\Mod(V)$ is independent of the choice of local density
or renormalization scheme.
\end{lemma}

\begin{proof}
Expand $e^{-f}\Delta_{\mathrm{cyc}}e^f$ on the bivector $P_\Pi$. Since $f$ has polyvector
degree zero, the new term is precisely the Hamiltonian bracket with $f$.
\end{proof}

\section{A vanishing theorem for triangular $W$-normal form}

We now identify a broad class of $W$-type PVAs for which the first variational modular
class vanishes identically.

\begin{definition}[Triangular $W$-normal form]
A filtered quadratic PVA is in \emph{triangular $W$-normal form} if it is strongly
generated by
\[
T,\ W^1,\dots,W^r
\]
such that:
\begin{enumerate}[leftmargin=2.4em]
\item $T$ is Virasoro:
\[
\{T_\lambda T\}=\partial T+2T\lambda+K(\lambda),
\]
with $K(\lambda)$ independent of the fields;
\item each $W^\alpha$ is primary of weight $h_\alpha$:
\[
\{T_\lambda W^\alpha\}=\partial W^\alpha+h_\alpha W^\alpha\lambda;
\]
\item in the row indexed by $W^\alpha$, the only dependence on jets of $W^\alpha$
appears in the entry $\Pi^{\alpha T}$.
\end{enumerate}
\end{definition}

\begin{theorem}[Triangular vanishing; \ClaimStatusProvedHere]\label{thm:triangular-vanishing}
If a filtered quadratic PVA is in triangular $W$-normal form, then
\[
\Delta_{\mathrm{cyc}}P_\Pi=(r+1)\int \partial\theta_T=0
\qquad\text{in }\Xloc^1/\partial.
\]
Consequently,
\[
\Mod(V)=0.
\]
\end{theorem}

\begin{proof}
The Virasoro row has operator
\[
\Pi^{TT}=T'+2T\partial+K(\partial),
\]
so
\[
\frac{\delta \Pi^{TT}}{\delta T}=(-\partial)\circ 1 + 2\partial=\partial.
\]

For a primary field $W^\alpha$ of weight $h_\alpha$,
\[
\Pi^{T\alpha}=(W^\alpha)'+h_\alpha W^\alpha\partial.
\]
Skew-adjointness gives
\[
\Pi^{\alpha T}=(h_\alpha-1)(W^\alpha)'+h_\alpha W^\alpha\partial,
\]
hence
\[
\frac{\delta \Pi^{\alpha T}}{\delta W^\alpha}
=-(h_\alpha-1)\partial+h_\alpha\partial=\partial.
\]
By assumption there are no further appearances of $W^\alpha$ in the $\alpha$-row. Therefore
the total modular operator contributes $\partial$ once from the Virasoro row and once from
each primary row:
\[
\Delta_{\mathrm{cyc}}P_\Pi=(r+1)\int \partial \theta_T.
\]
This vanishes modulo total derivatives.
\end{proof}

\begin{remark}
This theorem gives a large natural vanishing class. It says that in a broad swath of
classical $W$-type examples the first modular anomaly does \emph{not} appear as an
obstruction class. The quantum problem starts instead with the torsor of genus-one lifts.
\end{remark}

\section{The classical $W_3$ example}

We now specialize to the simplest nontrivial $W$-type example studied by Khan--Zeng,
namely the classical $W_3$ PVA.

\subsection{The classical $W_3$ brackets}

Let $T$ have spin $2$ and $W$ have spin $3$. The classical $W_3$ brackets are
\cite{KZ25}
\begin{align}
\{T_\lambda T\}_c
&=
\partial T+2T\lambda+\frac{c}{12}\lambda^3,
\\[4pt]
\{T_\lambda W\}_c
&=
\partial W+3W\lambda,
\\[4pt]
\{W_\lambda W\}_c
&=
\frac{c}{360}\lambda^5
+\frac13 T\lambda^3
+\frac12(\partial T)\lambda^2
+\Big(\frac3{10}\partial^2T+\frac{32}{5c}T^2\Big)\lambda
+\frac1{15}\partial^3T
+\frac{32}{5c}T\partial T.
\end{align}

The corresponding Hamiltonian operator is
\[
\Pi_{W_3}(\partial)=
\begin{pmatrix}
T'+2T\partial+\dfrac c{12}\partial^3
&
W'+3W\partial
\\[2mm]
2W'+3W\partial
&
\dfrac c{360}\partial^5
+\dfrac13 T\partial^3
+\dfrac12 T'\partial^2
+\Big(\dfrac{3}{10}T''+\dfrac{32}{5c}T^2\Big)\partial
+\dfrac1{15}T'''
+\dfrac{32}{5c}TT'
\end{pmatrix}.
\]

\begin{theorem}[Vanishing of the first modular shadow for $W_3$; \ClaimStatusProvedHere]\label{thm:W3mod0}
For the classical $W_3$ PVA,
\[
\Delta_{\mathrm{cyc}}P_{W_3}=2\int \partial\theta_T=0.
\]
Hence
\[
\Mod(W_3)=0.
\]
\end{theorem}

\begin{proof}
Compute the $T$-component of the variational modular operator:
\[
\mathfrak M_{W_3}^T
=
\frac{\delta \Pi^{TT}}{\delta T}
+
\frac{\delta \Pi^{WT}}{\delta W}.
\]
Now
\[
\frac{\delta \Pi^{TT}}{\delta T}=(-\partial)\circ 1+2\partial=\partial,
\]
while
\[
\Pi^{WT}=2W'+3W\partial
\qquad\Longrightarrow\qquad
\frac{\delta \Pi^{WT}}{\delta W}=(-\partial)\circ 2 + 3\partial=\partial.
\]
Therefore
\[
\mathfrak M_{W_3}^T=2\partial.
\]

For the $W$-component,
\[
\mathfrak M_{W_3}^W
=
\frac{\delta \Pi^{TW}}{\delta T}
+
\frac{\delta \Pi^{WW}}{\delta W}.
\]
But $\Pi^{TW}=W'+3W\partial$ has no $T$-dependence and $\Pi^{WW}$ has no $W$-dependence.
Hence $\mathfrak M_{W_3}^W=0$.

Thus
\[
\Delta_{\mathrm{cyc}}P_{W_3}=\int 2\partial\theta_T=0
\]
modulo total derivatives.
\end{proof}

\subsection{Central-parameter derivative as a cocycle}

The next structural fact is completely general.

\begin{theorem}[Parameter-derivative cocycle; \ClaimStatusProvedHere]\label{thm:parameter-cocycle}
Let $P_c$ be a one-parameter family of local Poisson bivectors satisfying
\[
[P_c,P_c]=0
\qquad\text{for all }c.
\]
Then
\[
\Xi_c:=\partial_c P_c
\]
is $\delta_{P_c}$-closed:
\[
[P_c,\Xi_c]=0.
\]
Equivalently, $\Xi_c$ determines a first-order deformation class.
\end{theorem}

\begin{proof}
Differentiate the identity $[P_c,P_c]=0$ with respect to $c$:
\[
0=\partial_c[P_c,P_c]=2[P_c,\partial_cP_c].
\]
\end{proof}

For the $W_3$ family one finds
\[
\partial_c \Pi^{TT}_{W_3}=\frac1{12}\partial^3,\qquad
\partial_c\Pi^{TW}_{W_3}=0,
\]
and
\[
\partial_c \Pi^{WW}_{W_3}
=
\frac1{360}\partial^5-\frac{32}{5c^2}(T^2\partial+T\partial T).
\]
Thus $\partial_c P_{W_3}$ is an explicit deformation cocycle.

\section{The relevant shifted $H^1$-sector for $W_3$}

We now compute the finite-dimensional local sector relevant to genus-one lifts.

\subsection{Definition of the relevant sector}

\begin{definition}[Relevant shifted deformation sector]
Let $P_c$ be the classical $W_3$ Poisson bivector. Define
\[
C^1_{\mathrm{rel}}(W_3)
\]
to be the space of local, translation-invariant, filtration-preserving, reduced conformal-weight
zero bivectors for the strong generators $(T,W)$. Equivalently, it is the finite-dimensional
ansatz of first-order deformations preserving the strong generating type and the conformal
degrees of the classical brackets.
\end{definition}

\begin{remark}[Shift convention]
Although the representatives in $C^1_{\mathrm{rel}}(W_3)$ are ordinary bivectors, we call
the resulting cohomology group ``shifted $H^1$'' because this is the degree in the shifted
Maurer--Cartan deformation complex. In ordinary Poisson cohomology this corresponds to
degree $2$.
\end{remark}

\subsection{Normal form of relevant cochains}

\begin{proposition}[Normal form; \ClaimStatusProvedHere]\label{prop:W3normalform}
Every element $Q\in C^1_{\mathrm{rel}}(W_3)$ has a unique expression
\begin{align}
Q_{TT}&=a\,\lambda^3+b\,(2T\lambda+\partial T),\label{eq:QTT}\\
Q_{TW}&=p\,\lambda^4+q\,T\lambda^2+(r\,\partial T+s\,W)\lambda+u\,\partial^2T+v\,T^2+w\,\partial W,\label{eq:QTW}\\
Q_{WW}&=
A\,\lambda^5
+B\Big(T\lambda^3+\frac32\partial T\,\lambda^2+\frac12\partial^2T\,\lambda\Big)
+C\,(2\partial^2T\,\lambda+\partial^3T)\notag\\
&\qquad
+D\,(T^2\lambda+T\partial T)
+E\,(2\partial W\,\lambda+\partial^2W),\label{eq:QWW}
\end{align}
with $Q_{WT}$ determined by skew-symmetry.
\end{proposition}

\begin{proof}
The reduced conformal-weight zero, locality, filtration-preserving, and skew-symmetry
constraints leave exactly the monomials displayed above.
\end{proof}

\subsection{Linearized cocycles}

Let
\[
\delta_{P_c}Q:=[P_c,Q]
\]
be the linearized differential.

\begin{proposition}[Linearized cocycle equations; \ClaimStatusProvedHere]\label{prop:W3cocycles}
A relevant cochain $Q$ of the form \eqref{eq:QTT}--\eqref{eq:QWW} is a cocycle if and only if
its coefficients satisfy
\begin{align}
r&=u=v=0,\qquad s=3b,\qquad w=b,\label{eq:rel1}\\
p&=\frac c{24}E,\qquad q=E,\label{eq:rel2}\\
B&=5C,\label{eq:rel3}\\
a&=\frac c{12}b+\frac{5c}{4}C-\frac{5c^2}{384}D,\label{eq:rel4}\\
A&=\frac c{12}C-\frac{c^2}{2304}D.\label{eq:rel5}
\end{align}
Hence the cocycle space $Z^1_{\mathrm{rel}}(W_3)$ is $4$-dimensional with free parameters
\[
b,\ C,\ D,\ E.
\]
\end{proposition}

\begin{proof}
One expands the linearized Jacobi identities for the triples $(T,T,W)$, $(T,W,W)$, and
$(W,W,W)$, then equates coefficients of the independent monomials in $\lambda$ and $\mu$.
This is a finite but lengthy coefficient computation; the relations above are the complete
solution.
\end{proof}

\begin{remark}[Explicit parametrization of cocycles]
Equations \eqref{eq:rel1}--\eqref{eq:rel5} give
\begin{align*}
Q_{TT}
&=
\Big(\frac c{12}b+\frac{5c}{4}C-\frac{5c^2}{384}D\Big)\lambda^3+b(2T\lambda+\partial T),\\[4pt]
Q_{TW}
&=
\frac c{24}E\,\lambda^4+E\,T\lambda^2+3b\,W\lambda+b\,\partial W,\\[4pt]
Q_{WW}
&=
\Big(\frac c{12}C-\frac{c^2}{2304}D\Big)\lambda^5
+5C\Big(T\lambda^3+\frac32\partial T\,\lambda^2+\frac12\partial^2T\,\lambda\Big)\\
&\qquad
+C(2\partial^2T\,\lambda+\partial^3T)
+D(T^2\lambda+T\partial T)
+E(2\partial W\,\lambda+\partial^2W).
\end{align*}
\end{remark}

\subsection{Exact directions}

\begin{proposition}[Relevant exact directions; \ClaimStatusProvedHere]\label{prop:W3exacts}
The reduced conformal-weight zero, local, translation-invariant evolutionary vector fields
are exactly
\[
X=x\,T\frac{\delta}{\delta T}+(y\,W+z\,\partial T)\frac{\delta}{\delta W}.
\]
Their Lie derivatives on the classical $W_3$ bivector are
\begin{align}
(L_XP_c)_{TT}
&=
-x\Big(\frac c6\lambda^3+2T\lambda+\partial T\Big),\label{eq:LXTT}\\
(L_XP_c)_{TW}
&=
-z\Big(\frac c{12}\lambda^4+2T\lambda^2\Big)-x(3W\lambda+\partial W),\label{eq:LXTW}\\
(L_XP_c)_{WW}
&=
-\frac{cy}{180}\lambda^5
+\Big(\frac x3-\frac{2y}{3}\Big)
\Big(T\lambda^3+\frac32\partial T\,\lambda^2+\frac12\partial^2T\,\lambda\Big)\notag\\
&\qquad
+\Big(\frac x{15}-\frac{2y}{15}\Big)(2\partial^2T\,\lambda+\partial^3T)
+\frac{64(x-y)}{5c}(T^2\lambda+T\partial T)
-2z(2\partial W\,\lambda+\partial^2W).\label{eq:LXWW}
\end{align}
Consequently the exact subspace $B^1_{\mathrm{rel}}(W_3)$ is $3$-dimensional.
\end{proposition}

\begin{proof}
The stated vector fields are exactly the local translation-invariant evolutionary fields
compatible with the filtration and reduced conformal weight. Applying the Lie derivative
to the three bracket entries gives \eqref{eq:LXTT}--\eqref{eq:LXWW}.
\end{proof}

\subsection{The cohomology computation}

\begin{theorem}[The relevant shifted $H^1$ for $W_3$; \ClaimStatusProvedHere]\label{thm:H1W3}
Assume $c\neq 0$. Then:
\begin{enumerate}[leftmargin=2.4em]
\item every cocycle $Q\in Z^1_{\mathrm{rel}}(W_3)$ decomposes uniquely as
\[
Q=L_XP_c+\kappa\,\partial_c P_c
\]
for some relevant vector field $X$ and some scalar $\kappa\in \mathbb C$;
\item the class of $\partial_c P_c$ is nonzero;
\item therefore
\[
H^1_{\mathrm{rel}}(W_3):=
Z^1_{\mathrm{rel}}(W_3)/B^1_{\mathrm{rel}}(W_3)
\cong \mathbb C\,[\partial_cP_c].
\]
\end{enumerate}
\end{theorem}

\begin{proof}
Let $Q$ be a cocycle parametrized by $(b,C,D,E)$ as in \cref{prop:W3cocycles}.
Define
\[
x:=-b,\qquad
y:=-\frac{15C+b}{2},\qquad
z:=-\frac E2,
\]
and
\[
\kappa:=c(15C-b)-\frac{5c^2}{32}D.
\]
A direct substitution into \eqref{eq:LXTT}--\eqref{eq:LXWW} shows that
\[
Q=L_XP_c+\kappa\,\partial_cP_c.
\]
Thus every cohomology class is represented by a scalar multiple of $\partial_cP_c$.

It remains to prove that $\partial_cP_c$ is not exact. Suppose
\[
\partial_cP_c=L_XP_c
\]
for some relevant vector field $X$. From the $TW$-component one immediately obtains
$z=0$. Matching the central term in the $TT$-component forces $x=-1/(2c)$, while the
central term in the $WW$-component forces $y=-1/(2c)$. But then the coefficient of
$T^2\lambda+T\partial T$ in $L_XP_c$ is
\[
\frac{64(x-y)}{5c}=0,
\]
whereas in $\partial_cP_c$ it equals
\[
-\frac{32}{5c^2}\neq 0.
\]
This contradiction proves non-exactness.
\end{proof}

\begin{corollary}[Uniqueness of the relevant genus-one lift direction; \ClaimStatusProvedHere]\label{cor:unique-lift}
In the local, translation-invariant, filtration-preserving, reduced conformal-weight zero
sector, every first-order genus-one lift of the classical $W_3$ structure is gauge equivalent
to a unique scalar multiple of the central-parameter direction:
\[
P_c\longmapsto P_c+\hbar\,\kappa\,\partial_cP_c
=
P_{c+\hbar\kappa}\pmod{\hbar^2}.
\]
Equivalently, the relevant shifted $H^1$ is one-dimensional and generated by
$[\partial_cP_c]$.
\end{corollary}

\begin{proof}
This is the content of \cref{thm:H1W3}.
\end{proof}

\begin{remark}[Why this is the correct rigidity statement]
We have not computed the full local Poisson cohomology of $W_3$ in all degrees and all
filtration sectors. What we have computed is precisely the sector that matches the locality,
translation invariance, filtration, and conformal behavior expected of the one-loop boundary
problem in the Khan--Zeng model. In that physically and algebraically natural sector, the
central shift is the only deformation parameter.
\end{remark}

\section{Synthesis: the modular quantization program}

We now synthesize the entire picture.

\subsection{Disk level}

At genus zero, the deformation problem is controlled by the ordinary tree-level chiral bar
complex. In the dualizable quadratic setting this recovers the Gui--Li--Zeng Maurer--Cartan
picture
\[
\Hom(A,B)\hookrightarrow \MC(A^!\otimes B),
\]
and in the inhomogeneous quadratic setting it recovers the twisted equation
\[
\mu\big((S+\alpha)\boxtimes(S+\alpha)\big)=0.
\]
This is the classical, disk-level, genus-zero shadow.

\subsection{The modular correction}

The new modular input is the nonseparating one-edge operator $D_1$. On the coordinate side
it becomes the reduced cyclic Laplacian $\Delta_{\mathrm{cyc}}$. The first forcing term in the
modular Maurer--Cartan recursion is the class of $d_1\Theta_0$; its coordinate shadow is the
variational modular class
\[
\Mod(V)=[\Delta_{\mathrm{cyc}}P_\Pi].
\]

\subsection{The $W$-algebra lesson}

For triangular $W$-normal form the first modular shadow vanishes identically. In particular,
for classical $W_3$,
\[
\Mod(W_3)=0.
\]
Thus the genus-one problem begins not with an obstruction but with a torsor of lifts. The
computation of \cref{thm:H1W3} shows that in the relevant sector this torsor is one-dimensional
and generated by $\partial_cP_c$.

\subsection{The conditional physical interpretation}

Khan--Zeng predict that quantization of the classical $W_3$ PVA leads to the quantum $W_3$
algebra with a shifted effective central parameter \cite{KZ25}. The formal algebraic
results of the present note explain why such a shift is the \emph{only available relevant
genus-one direction} once locality, conformal weight, and filtration are imposed. What
remains is the genuinely analytic and physical problem of computing the scalar $\kappa$
selected by the three-dimensional Poisson sigma model. The literature points toward
$\kappa=50$ in the $W_3$ case, but this numerical identification is not derived here.

\section{Outlook and frontier problems}

The constructions and theorems above open several precise frontier problems.

\begin{question}[Actual extraction of the one-edge operator from the 3d sigma model]
Construct the nonseparating one-edge operator $D_1$ directly from the one-loop Feynman
graphs of the Khan--Zeng Poisson sigma model, and prove that its boundary shadow is
$\Delta_{\mathrm{cyc}}$.
\end{question}

\begin{question}[The scalar $\kappa$]
Compute the coefficient $\kappa$ of the unique relevant genus-one lift direction for
$W_3$ from the half-space theory itself. This would convert the qualitative rigidity
result of \cref{cor:unique-lift} into an actual one-loop quantization theorem.
\end{question}

\begin{question}[Beyond $W_3$]
Extend the triangular vanishing theorem and the relevant shifted $H^1$ computation to
higher classical $W_N$-algebras and to Drinfeld--Sokolov reductions of general type.
\end{question}

\begin{question}[Curved modular bar/cobar duality]
Construct the full curved modular bar--cobar adjunction for twisted pairs $(B,B^\circ,S)$.
The point is not merely to add curvature to formulas, but to preserve the genuinely
extra information contained in the twisted-pair data of Gui--Li--Zeng.
\end{question}

\begin{question}[Line operators and modular chiral quantization]
Dimofte--Niu--Py show that line operators in perturbative 3d holomorphic QFT are governed
by a Koszul-dual $A_\infty$ algebra $A^!$ together with Maurer--Cartan data of $r$-matrix
type \cite{DNP25}. It is natural to ask for a common modular formalism that
simultaneously governs:
\begin{enumerate}[leftmargin=2.4em]
\item bulk local operators,
\item boundary chiral algebras,
\item line-operator OPE,
\item and the genus filtration coming from stable gluing.
\end{enumerate}
The present note suggests that the stable-graph modular deformation complex should be the
right place to start.
\end{question}

\section*{Epilogue}

The disk knows the classical bracket; the loop knows the modular correction. Trees encode
chirality, loops encode modularity, and the stable graph is the bridge between them.
What the present note adds is a concrete algebraic bridge:
\[
\text{PVA}
\rightsquigarrow
\text{resolved chiral datum}
\rightsquigarrow
\Bmod
\rightsquigarrow
\Defmod
\rightsquigarrow
\text{modular Maurer--Cartan theory}.
\]
At genus zero this recovers the known quadratic-dual picture. At genus one it produces
a concrete new invariant, the variational modular class. For $W_3$ that invariant vanishes,
and the remaining lift freedom is one-dimensional, generated by the central-parameter
direction. That is the precise point at which the analytic problem of quantization can
now engage the algebraic problem of modular deformation theory.

